\documentclass[fleqn]{article}
%\textwidth=6in
%\textheight=9in
%\pagestyle{empty}
\usepackage{amsmath,amssymb,amsthm,ragged2e,siunitx,graphicx,multicol,wrapfig,mathtools,parskip}
\usepackage[margin=0.7in]{geometry}
\usepackage[T1]{fontenc}
\usepackage[utf8]{inputenc}
\usepackage[dvipsnames]{xcolor}
\usepackage[export]{adjustbox}
\newtheorem{thm}{{\sc Theorem}}[section]
\newtheorem{prop}[thm]{{\sc Proposition}}
\newtheorem{cor}[thm]{{\sc Corollary}}
\newtheorem{lem}[thm]{{\sc Lemma}}
\newtheorem{df}[thm]{{\sc Definition}}

\DeclarePairedDelimiter\ceil{\lceil}{\rceil}
\DeclarePairedDelimiter\floor{\lfloor}{\rfloor}

\DeclareMathOperator{\sech}{sech}
\DeclareMathOperator{\csch}{csch}

% automatically makes new paragraphs have no indent
\setlength{\parindent}{0pt}

\newcommand{\n}{\noindent}
\newcommand{\ds}{\displaystyle}
\newcommand{\rar}{\rightarrow}

% new command for underscore
\newcommand{\underscore}{\underline{\hspace{1cm}}}

% this is a command for the space between asterisk or plus and the problem statement
\newcommand{\gap}{\hspace{.25cm}}

\newcommand*\colvec[1]{\begin{bmatrix*}[r]#1\end{bmatrix*}} 

% Uncomment this to write vectors bolded 
\renewcommand{\vec}[1]{\boldsymbol{#1}}

% Uncomment this to write vectors bolded AND with arrow on top
\newcommand{\avec}[1]{\overrightarrow{\boldsymbol{#1}}}
% this notation has an arrow on top of vector which is redundant if it is bolded

% Uncomment this to write vectors unbolded with long arrow on top
% \renewcommand{\vec}[1]{\overrightarrow{#1}}

% this is a command for a coefficient matrix with right-aligned entries (good for aligning negative signs)
\newcommand{\mat}[1]{\begin{bmatrix*}[r]#1\end{bmatrix*}} 

% this is a command for a coefficient matrix with center-aligned entries (good for aligning negative signs)
\newcommand{\cmat}[1]{\begin{bmatrix*}[c]#1\end{bmatrix*}} 

% this is a command for a DETERMNINANT matrix with right-aligned entries (good for aligning negative signs) (ABSOLUTE VALUE BARS)
\newcommand{\vmat}[1]{\begin{vmatrix*}[r]#1\end{vmatrix*}} 

% this is a command for black squares (useful for pivots in matrices) when already in math mode (between $ signs)
\newcommand{\bsq}{\blacksquare}






%%  environment for AUGMENTED  matrices
\newenvironment{amatrix}[1]{%
  \left[\begin{array}{@{}*{#1}{r}|r@{}}
}{%
  \end{array}\right]
}


%%% new command for Augmented BAR alone (helpful for putting augmented bar somewhere in middle of matrix)
\newcommand\aug{\fboxsep=-\fboxrule\!\!\!\fbox{\strut}\!\!\!}

% define new color (i.e. pink)
\definecolor{mypink1}{rgb}{0.858, 0.188, 0.478}

% define new color (i.e. medium green)
\definecolor{MediumGreen}{rgb}{0.1, 0.8, 0.65}

% define new color (i.e. dark purple)
\definecolor{DarkPurple}{rgb}{0.6, 0, 0.6}

% this is a command for the size of headers
\newcommand{\header}[1]{\noindent{\Large \textcolor{mypink1}{\textbf{#1}}}}

% this is a command for the size of subheaders
\newcommand{\subheader}[1]{\noindent{\large \textcolor{red}{\textbf{#1}}}}

% this is a command for answers in RED on same line with a space between problem and answer
\newcommand{\sred}[1]{\hspace{1cm} \textcolor{red}{#1}}

% this is a command for answers in RED on same line with NO space between problem and answer
\newcommand{\red}[1]{\textcolor{red}{#1}}

% this is a command for answers in PURPLE on same line with NO space between problem and answer
\newcommand{\DarkPurple}[1]{\textcolor{DarkPurple}{#1}}

% this is a command for answers in GREEN on same line with NO space between problem and answer
\newcommand{\MediumGreen}[1]{\textcolor{MediumGreen}{#1}}

% this is a command for answers in BLUE on same line with a space between problem and answer
\newcommand{\sblue}[1]{\hspace{1cm} \textcolor{blue}{#1}}

% this is a command for answers in BLUE on same line with NO space between problem and answer
\newcommand{\blue}[1]{\textcolor{blue}{#1}}


% this is a command for defining EXAMPLES
\theoremstyle{definition}
\newtheorem{exmp}{Example}[section]




\begin{document}

\centerline{\Large \bf Math Seminar Template Notes}
\centerline{\large \bf Day 6 (Proof by Contraposition and Proof by Contradiction)}
\vspace{\baselineskip}

%\noindent \underline{\hspace{18cm}}

\header{Day 6 (Proof by Contraposition and Proof by Contradiction) } \\


\noindent\fbox{%
\parbox{\textwidth}{
\subheader{Overview} \\ 
%\vspace{0.25cm}
\textbf{\underline{Goals}} 
\begin{enumerate}
\item Do examples using \textbf{proofs by contraposition}
\item Introduce and apply a \textbf{proof by contradiction}
\end{enumerate}
\vspace{.25cm}
}}

\vspace{0.2cm}

\subheader{Review of Proof by Contraposition}

%%%% something in a box
\noindent\fbox{%
\parbox{\textwidth}{\vspace{.25cm}
\textbf{Def:} If we have a statement of the form ``if $p$, then $q$" (symbolically written $p \rightarrow q$), then the \textbf{contrapositive} is

\vspace{.25cm}
\hspace{5cm} ``if not $q$, then not $p$" (symbolically written $\lnot q \rightarrow \lnot p$). 

\vspace{.5cm}
These two statements are \textbf{logically equivalent}.


\vspace{.25cm}   
}}
%%% end box

\vspace{.25cm}

\textbf{Ex 1:} Write the contrapositives of the following statements: 

\vspace{.25cm}
(a) If it's raining, then the grass is wet.


\vspace{1.5cm}


% (b) If $n^2$ is even, then $n$ is even. 

% \vspace{1.5cm}
(b) On Wednesdays, we wear pink. 

\vspace{1.5cm}


% \vspace{.25cm}
%%%% something in a box
\noindent\fbox{%
\parbox{\textwidth}{\vspace{.25cm}
If we are trying to prove the statement ``if $p$, then $q$" using a \textbf{proof by contraposition}, we instead try to show that the contrapositive, ``if not $q$, then not $p$" is true by doing the following:

    \vspace{.25cm}
    \begin{itemize}
    \item We assume $q$ does not happen (i.e. $q$ is false)

    \vspace{.1cm}
    \item We try to show that this means as a result $p$ will not happen (i.e. $p$ is false)
    \end{itemize}
\vspace{.25cm}   
}}
%%% end box

\vspace{.25cm}

\textbf{Ex 2 (last time):} Let's revisit the proposition we saw previously:

\vspace{.25cm}
    \hspace{1cm} \textbf{Proposition:} If $n^2$ is even, then $n$ is even.

\vspace{.25cm}
(a) Write the contrapositive of this statement. 

\vspace{1.5cm}
(b) Why is it hard to do a direct proof here?


%%%
\newpage


(c) Use a proof by contraposition to prove the statement. 

%%% Scratch work for proof
\hspace{7cm} \MediumGreen{\underline{Scratch Work}}
\begin{multicols}{2} 
\hspace{4cm} \MediumGreen{\underline{Given/Know}} \columnbreak

\MediumGreen{\underline{Want to show}} 
\end{multicols}

\vspace{6cm} 

\textbf{Pf:}



%%%
\newpage

\textbf{Ex 3:} Consider the following proposition: \hspace{.25cm} For an integer $n$, if $3n+5$ is even, then $n$ is odd. 
%Show the sum of two odd integers is even.

\vspace{.5cm}
(a) First try proving this using a direct proof. What difficulty do you run into? 

%%% Scratch work for proof
\hspace{7cm} \MediumGreen{\underline{Scratch Work}}
\begin{multicols}{2} 
\hspace{4cm} \MediumGreen{\underline{Given/Know}} \columnbreak

\MediumGreen{\underline{Want to show}} 
\end{multicols}

\vspace{4cm}
(b) Now try proving it using a proof by contraposition. 

\vspace{1cm}
%%% Scratch work for proof
\hspace{7cm} \MediumGreen{\underline{Scratch Work}}
\begin{multicols}{2} 
\hspace{4cm} \MediumGreen{\underline{Given/Know}} \columnbreak

\MediumGreen{\underline{Want to show}} 
\end{multicols}

\vspace{5cm}


%%%
\newpage


\subheader{Motivation for Proof by Contradiction}

%%%% something in a box
\noindent\fbox{%
\parbox{\textwidth}{\vspace{.25cm}
\textbf{Def:} A \textbf{rational} number is a real number that can be written in the form $\ds{\frac{a}{b}}$ where $a, b$ are integers and $b$ \underline{\hspace{1cm}}

\vspace{.5cm}
\hspace{1cm} \textbf{Examples:}

\vspace{1cm}
\textbf{Def:} An \textbf{irrational} number is a real number that is NOT rational. 

\vspace{.5cm}
\hspace{1cm} \textbf{Examples:}

\vspace{.5cm}   
}}
%%% end box

\vspace{.25cm}
\textbf{Ex 4:} Consider the proposition, ``the sum of a rational number and an irrational number is irrational." 

\vspace{.25cm}
First, let's restate this formally using variables. 

\vspace{.25cm}
    \hspace{1cm} ``If $x$ is \underline{\hspace{3cm}} and $y$ is \underline{\hspace{3cm}}, then \underline{\hspace{1.5cm}} is \underline{\hspace{3cm}}"

\vspace{.5cm}
(a) What makes doing a direct proof difficult? 

\vspace{2cm}
(b) What is the contrapositive of this statement? What makes doing a proof by contraposition difficulty?

\vspace{4cm}
\textbf{Ex 5:} Consider the following scenario. Suppose that everyone in our classroom is here and there's also a newborn baby in the room (idk why!). I notice that the camera at the back of the classroom has dismantled. 

\vspace{.25cm}
I decide to do a thought experiment. I assume that the baby did it. 

\vspace{.25cm}
(a) What makes my assumption kind of ridiculous? 

\vspace{1.5cm}
(b) Because my assumption is \underline{\hspace{3cm}}, what does this tell us? 

\vspace{1.5cm}


%%%
\newpage 

%%%% something in a box
\noindent\fbox{%
\parbox{\textwidth}{\vspace{.25cm}
In a \textbf{proof by contradiction}

    \vspace{.25cm}
	\hspace{1cm} (a) Assume that what we want to prove, is \red{false}

	\vspace{.25cm}
	\hspace{1cm} (b) Show this leads to a \red{contradiction} (e.g. ``integer $= \displaystyle{\frac{1}{2}}$," \hspace{.5cm} ``rational $= \underline{\hspace{3cm}},$ etc.

	\vspace{.25cm}
	\hspace{1cm} (c) This means our assumption is \underline{\hspace{3cm}} 

	\vspace{.25cm}
	\hspace{1.5cm} This implies that we wanted to prove is \underline{\hspace{2.5cm}}
\vspace{.25cm}   
}}
%%% end box

\vspace{.25cm}

\textbf{Ex 6:} Let's return to our earlier proposition: ``the sum of a rational number and an irrational number is irrational." 

\vspace{.25cm}
    \hspace{1cm} We stated this formally as ``if $x$ is rational and $y$ is irrational, then $x+y$ is irrational" 

\vspace{.5cm}
Let's try proving this using a \textbf{proof by contradiction}.

\vspace{1cm}
We will assume that \underline{\hspace{7cm}} and include this as one of our givens. 

%%% Scratch work for proof
\hspace{7cm} \MediumGreen{\underline{Scratch Work}}
\begin{multicols}{2} 
\hspace{4cm} \MediumGreen{\underline{Given/Know}} \columnbreak

\MediumGreen{\underline{Want to show}} 
\end{multicols}

\vspace{5cm} 

\textbf{Pf:}






%%%
\newpage

\subheader{Practice Problems}


\vspace{.25cm}


\begin{enumerate}

\item Consider the following proposition: 

\vspace{.25cm}
    \hspace{1cm} \textbf{Proposition:} If $n^2$ is odd, then $n$ is odd.

\vspace{.25cm}
Prove this using a proof by contradiction. Is it different from a proof by contraposition?



%%%
\newpage


\item Prove that product of two rational numbers is rational. 


%%%
\newpage


\item Prove that the product of a nonzero rational number and an irrational number is irrational. 






\end{enumerate}




\end{document}








\subheader{What Can We Assume is True and What We Cannot Assume is True?}

\vspace{.25cm}
%%%% something in a box
\noindent\fbox{%
\parbox{\textwidth}{\vspace{.25cm}
\underline{Things we can assume are true} 

\vspace{.25cm}

	\begin{enumerate}
	\item Axioms of real numbers and integers. 

		\begin{itemize}
		\item If $x, y$ are real numbers then $x \pm y$, $xy$, and $\displaystyle{\frac{x}{y}}$ are all real (except if $y=0$ in the division) 

		\vspace{.25cm}

		\item If $x, y$ are integers, then $x \pm y$ and $xy$ are integers. 

		\vspace{.25cm}

		\item \textbf{Distributive Law:} If $x, y, z$ are reals, then $x(y+z) = xy+xz$  

		\vspace{.25cm}

		\item \textbf{Commutative Law:} If $x, y$ are reals, then $x+y=y+x$ and $xy=yx$ 

		\vspace{.25cm}

		\item \textbf{Associative Law:} If $x, y, z$ are reals, then $x+(y+z) = (x+y)+z$. 

		\vspace{.25cm}

		\item \textbf{Transitive Law:} If $x, y, z$ are reals with $x<y$ and $y<z$, then $x<z$. (similarly for other inequalities)
		\end{itemize}

	\vspace{.25cm}
	\item Basic Arithmetic Rules 

		\begin{itemize}
		\item Rules for adding/subtracting/multiplying/dividing fractions (e.g. $\displaystyle{ \frac{a}{b} \cdot \frac{c}{d} = \frac{ac}{bd} }$ ) 

		\vspace{.25cm}

		\item Multiplication by $-1$: e.g. if $a$ is rational, then $-a$ is rational (or if $a$ is an integer, then $-a$ is an integer, etc.) 

		\vspace{.25cm}

		\item Multiplication by $0$: If $a$ is real, then $a \cdot 0 = 0$ 

		\vspace{.25cm}

		\item Let $x, y$ be real. If $xy = 0$, then $x=0$ or $y=0$
		\end{itemize}
	\end{enumerate}
\vspace{.25cm}   
}}
%%% end box


\vspace{.5cm}

%%%% something in a box
\noindent\fbox{%
\parbox{\textwidth}{\vspace{.25cm}
\red{\underline{Examples of things we cannot assume are true}} 

\vspace{.25cm}

	\begin{enumerate}
	\item ``Odd + Odd = Even" \hspace{.5cm} (or other versions of this like ``Odd + Even = Odd") 	

	\vspace{.25cm}

	\item Rational + rational = rational 

	\vspace{.25cm}

	\item Generally any statement more complicated than the ones in the axioms of real numbers and integers 
	\end{enumerate}
\vspace{.25cm}   
}}
%%% end box










%%%% something in a box
\noindent\fbox{%
\parbox{\textwidth}{\vspace{.25cm}
\textbf{Strategies to Show Inductive Step}

\vspace{.25cm}
	\begin{enumerate}
	\item Start with $P(k)$ and do stuff to \underline{\hspace{4cm}} so it becomes $P(k+1)$

	\vspace{.25cm}
	\item Start with one side/part of $P(k+1)$ and use the assumption that \underline{\hspace{4cm}} to get the rest of $P(k+1)$
	\end{enumerate}
\vspace{.25cm}   
}}
%%% end box




\begin{multicols}{5} 
Type 
\columnbreak 

Symbol 
\columnbreak

Definition 
\columnbreak

\hspace{1cm}
 \columnbreak
 
\hspace{1cm}
\end{multicols}






%%%%%%%%%%%%%%%%%%%%%%%%%%%%%%%%%%%%%%%%%%%%%%%%%%%%%%%%%




%image aligned
%%
\includegraphics[scale=0.7, right]{BlankAxes}
\vspace{-10pt}


%image centered
%%
\begin{center}
\includegraphics[scale=1.4]{SinCosTanDrawGraphs}
\end{center} 
\vspace{-20pt}



%%% incorrect/correct show something with 2 approaches
\vspace{.25cm} 
\begin{multicols}{2} 
\red{\underline{Incorrect}} \columnbreak

\MediumGreen{\underline{Correct}} 
\end{multicols}


%%% Scratch work for proof
\hspace{7cm} \MediumGreen{\underline{Scratch Work}}
\begin{multicols}{2} 
\hspace{4cm} \MediumGreen{\underline{Given/Know}} \columnbreak

\MediumGreen{\underline{Want to show}} 
\end{multicols}

\vspace{2.5cm} 



%%%% something in a box
\noindent\fbox{%
\parbox{\textwidth}{\vspace{.25cm}
\textbf{Def:} \\ (1) The \DarkPurple{intersection} of sets $A$ and $B$ is the set of all elements in \textbf{both} $A$ \textbf{and} $B$, denoted $\underline{\hspace{2cm}}$. 
\vspace{.25cm}   
}}
%%% end box


%%% big row reduction arrow
\xrightarrow{\hspace*{2cm}}  



%%% table with extra vertical spacing
\begin{tabular}{ c|c|c } 
$f'(x)$ &  $f(x)$  &  $F(x)$  \\ \hline
\rule{0pt}{3ex}           &  $x^n$  &             \\
 \rule{0pt}{5ex}          &  $\displaystyle{x^{-1}=\frac{1}{x}}$   &             \\
\rule{0pt}{5ex}           &  $e^x$  &             \\
\rule{0pt}{5ex}           &  $e^{kx}$  &             \\
\rule{0pt}{5ex}           &  $g(x) \pm h(x)$  &             \\
\rule{0pt}{5ex}           &  $cg(x)$  &             \\
\rule{0pt}{5ex}           &  $g(x)h(x)$  &             \\
\rule{0pt}{5ex}           &  $\displaystyle{\frac{g(x)}{h(x)}}$  &             \\
\end{tabular}



%%% coefficient matrix template
$A = \mat{ 1 & 0 & 1 & 2 & 2 \\ 0 & 0 & 1 & 0 & 1 \\ 2 & 0 & 1 & 4 & 3 }$


%%% augmented matrix template (the number in brackets indicates which column the augmented bar will be placed AFTER)
$\begin{amatrix}{5}
1 & -1 & 0 & -5 & 0 & -9  \\ 
0 & 0  & 1 &  4 & 0 & 4 \\ 
0 & 0  & 0 &  0 & 1 & 3  
 \end{amatrix}$
 
 
 
 %%%% template for multiple alignment points in system of equations
 $\begin{alignedat}{3}
		x_1 &+ x_2 &&- x_3 &&= -6 \\
		3x_1 &+ x_2 &&+x_3  &&= 8 \\
		4x_1 &+ 2x_2 &&      &&= 2 \\
		\end{alignedat}$
